\chapter[Finite arithmetic image and factorization \(468=18\cdot26\)]{Finite arithmetic image and factorization \(468=18\cdot26\)}
\label{chap:alfabeto-42}
\index{admissible decimal code}
\index{check digit}

The image in base \(1000\) of the finite map that assigns the decimal blocks is a
proper subset of the one thousand decimal blocks. The evolution law defined in the
\cref{chap:tpk-definicion} and the arithmetic of \(\mathbb Z/9\mathbb Z\)
determine an image of 42 triples, organized into six classes of seven. The
same separation of the additive and multiplicative observables subsequently factors the 468
outputs into 18 additive signatures and 26 multiplicative signatures.

\section{Primitives and recurrence of the finite sender}

\begin{definition}[Initial data of \(T_{\min}\)]
\label{pf84:E02-15}
Let \(D=\{N,E,S,O\}\) be given, with cardinal vectors
\[
 v(N)=(-1,0),\quad v(E)=(0,1),\quad
 v(S)=(1,0),\quad v(O)=(0,-1)
 \qquad\text{in }(\mathbb Z/9\mathbb Z)^2.
\]
The temporal signature and the multiplicative index are fixed by
\[
 \operatorname{sgn}_9(t)=
 \begin{cases}
 +1,&t\bmod9\in\{1,4,7\},\\
 -1,&t\bmod9\in\{2,5,8\},\\
 0,&t\bmod9\in\{0,3,6\},
 \end{cases}
\]
and
\[
 \operatorname{ind}_2(1,2,4,8,7,5)=(0,1,2,3,4,5),
\]
with value zero outside the units. An entry of \(T_{\min}\) is a pair
of seeds
\[
 I_9:=\{0,1,\ldots,8\},\qquad
 s_+=(i_+,j_+,d_+),\qquad s_\times=(i_\times,j_\times,d_\times)
 \in I_9^2\times D.
\]
These tables, the readers \(L_\Sigma,L_\Pi\), the six blocks of nine
instants, and the inversion of both vectors when \(27\mid t\) are all the
starting data. Neither \(\pi\), \(e\), \(\varphi\), nor any other
target value is introduced.
\end{definition}

\begin{algorithm}[54-step recurrence]
\label{pf84:E02-16}
At the outset, the two cursors occupy their seed positions and carry
the vectors \(v(d_+)\), \(v(d_\times)\). For
\(t=1,\ldots,54\), in this order:
\begin{enumerate}
  \item if \(\operatorname{sgn}_9(t)=+1\), \(L_\Sigma(i_+(t),j_+(t))\) is accumulated
  and the additive cursor is advanced;
  \item if \(\operatorname{sgn}_9(t)=-1\), then
  \(\operatorname{ind}_2(L_\Pi(i_\times(t),j_\times(t)))\) is accumulated and the
  multiplicative cursor is shifted;
  \item if \(\operatorname{sgn}_9(t)=0\), both cursors remain stationary;
  \item after that update, both vectors are reversed if
  \(27\mid t\).
\end{enumerate}
The accumulators are reset at the beginning of each block
\(9(r-1)+1\leq t\leq9r\). If \(S_r,E_r\) are their final values, then there is emitted
\[
 U_r=100\bar S_r+10(\bar S_r+\bar E_r)_{10}
       +(3S_r+5E_r+1)_{10},
 \qquad
 \bar S_r=S_r\bmod10,\quad \bar E_r=E_r\bmod10,
\]
y \(w_r=(-U_r)\bmod3\). Thus it is obtained.
\[
 T_{\min}(s_+,s_\times)=(U_1,\ldots,U_6)
 \in\mathcal U_6^{\rm cat},
 \qquad
 \Psi(U_6)=(-U_6)\bmod3=:w_6\in\mathbb F_3^6.
\]
\end{algorithm}

\begin{criterion}[Target-value-independent dual control] \label{pf84:E02-17} The recurrence admits two independent formulations: channel factorization and joint simulation. Exhaustive traversal of the \(104\,976\) inputs produces the same sextuple by both routes for every seed, without consulting a target constant. A discrepancy for a single seed would refute the equivalence.
\end{criterion}

\section{Bivariate dynamics and block observables}

We work on the torus \(\mathbb Z_9^2\). Under the preceding convention, a
complete seed is
\[
\omega=(p_+,p_\times,h_+,h_\times)
\in(\mathbb Z_9^2)^2\times\{N,E,S,O\}^2,
\]
so that
\[
|\Omega|=9^4\cdot4^2=104\,976.
\]
In each window of nine steps, the times \(1,4,7\) evaluate the additive observable,
the times \(2,5,8\) evaluate the multiplicative observable, and the times \(3,6,9\)
increment the neutral counter. If \(S_m,E_m,T_m\) are the accumulators
of the window \(m\), the decimal observable is
\[
d_1\equiv S_m,\qquad
d_2\equiv S_m+E_m,\qquad
d_3\equiv3S_m+5E_m+7T_m\pmod{10},
\]
and the decimal triple is \(U_m=100d_1+10d_2+d_3\).

\begin{lemma}[Neutral counter]
For every window, \(T_m=3\).
\end{lemma}
\begin{proof}
Each block of nine times contains exactly the three neutral times
\(3,6,9\). Therefore the counter is incremented three times.
\end{proof}

\section{Image of the additive accumulator: \(\operatorname{Im}S_m=\{6,9,12,15,18,21,24\}\)}

Let \(s=i+j\pmod9\). The visible additive observable is
\[
f(s)=\operatorname{dr}_9(s+2)=(2,3,4,5,6,7,8,9,1).
\]
A cardinal translation modifies \(s\) in \(+1\) or in \(-1\). In a window,
three consecutive terms are read; reversing the orientation only reverses their
order. The nine cyclic sums are
\[
9,12,15,18,21,24,18,12,6.
\]

\begin{lemma}[Image of the additive accumulator]
\[
S_m\in\{6,9,12,15,18,21,24\}.
\]
In particular, \(d_1=S_m\bmod10\) traverses precisely
\(\{1,2,4,5,6,8,9\}\).
\end{lemma}
\nopagebreak[4]
\begin{proof}
The preceding list enumerates all initial positions of the trajectory and the
two orientations. Their decimal residues are seven and distinct.\qedhere
\end{proof}

\section{Image of the multiplicative accumulator: \(\operatorname{Im}E_m=\{0,1,3,5,7,9\}\)}

The unit group \((\mathbb Z/9\mathbb Z)^\times\) is cyclic of order
six. We fix the multiplicative coordinate function
\[
\ell_\times(1,2,3,4,5,6,7,8,9)=(0,1,0,2,5,0,4,3,0).
\]
If the row or column traversed has a factor divisible by three, each
evaluation belongs to \(\{3,6,9\}\) and contributes zero. If the fixed factor is a
unit, the sums of three consecutive terms of
\(b\mapsto\ell_\times(\operatorname{dr}_9(ab))\) are:
\[
\begin{array}{c|c}
a&\text{possible sum}\\
\midrule
1&\{1,3,7,9\}\\
2&\{3,5,9\}\\
4&\{1,5,7\}\\
5&\{1,5,7\}\\
7&\{3,5,9\}\\
8&\{1,3,7,9\}.
\end{array}
\]

\begin{lemma}[Image of the multiplicative accumulator]
\[
E_m\in\{0,1,3,5,7,9\}.
\]
\end{lemma}
\begin{proof}
The noninvertible case contributes \(0\). The finite table of the six units
contains all possibilities in the invertible case, and their union is
\(\{1,3,5,7,9\}\).
\end{proof}

\section{Decimal congruence of the ternary alphabet: uniqueness of the 3rd digit}

\begin{theorem}[Control congruence]
Every terna \(\overline{d_1d_2d_3}\) emitted by the system satisfies
\[
\boxed{d_3\equiv5d_2-2d_1+1\pmod{10}.}
\]
\end{theorem}
\begin{proof}
From \(d_2\equiv S_m+E_m\) one obtains
\(E_m\equiv d_2-d_1\pmod{10}\). Since \(T_m=3\),
\[
d_3\equiv3d_1+5(d_2-d_1)+21
\equiv5d_2-2d_1+1\pmod{10}.\qedhere
\]
\end{proof}

\begin{theorem}[Exact image of the decimal application]
The system produces exactly \(7\cdot6=42\) decimal triples. Grouped by
\(E=d_2-d_1\pmod{10}\), they are:
\[
\begin{array}{c|rrrrrrr}
E&\multicolumn{7}{c}{\text{decimal triples}}\\
\midrule
0&114&227&443&556&669&885&998\\
1&129&232&458&561&674&890&903\\
3&149&252&478&581&694&810&923\\
5&169&272&498&501&614&830&943\\
7&189&292&418&521&634&850&963\\
9&109&212&438&541&654&870&983.
\end{array}
\]
\end{theorem}
\begin{proof}
There are seven possibilities for \(d_1\) and six for \(E\). For each pair,
\(d_2\equiv d_1+E\) and the control congruence fix a unique \(d_3\),
so there exist at most 42 triples. The table applies these two formulas
to the 42 pairs and contains no repetitions. The enumeration of the
104\,976 seeds finds all the triples in the table and none outside
it; the bound is attained.
\end{proof}

\section{Factorization of the \(2\) observables}

Let \(\mathcal S=\mathbb Z_9^2\times\{N,E,S,O\}\), with
\(|\mathcal S|=324\). To each additive seed \(P\in\mathcal S\) is
associated the six-window signature
\[
s(P)=(S_1,\ldots,S_6)\pmod{10},
\]
and to each multiplicative seed \(T\in\mathcal S\),
\[
e(T)=(E_1,\ldots,E_6)\pmod{10}.
\]
The two signatures are combined coordinate by coordinate by
\[
G(s,e)_k=
100s_k+10(s_k+e_k)_{10}+(3s_k+5e_k+1)_{10}.
\]
The map \(G\) is injective: from a triple one recovers
\(s_k=d_1\) and \(e_k=d_2-d_1\pmod{10}\), while the third digit verifies
membership in the image.

\begin{lemma}[Exact classification of the additive channel]
\label{pf84:E02-18}
\[
 |\operatorname{im}s|=18,
\]
and each additive signature has exactly \(18\) preimages.
\end{lemma}
\begin{proof}
The signature depends only on
\((i+j\bmod9,\varepsilon)\), where \(\varepsilon=+1\) for \(E,S\) and
\(-1\) for \(N,O\). The \(9\cdot2\) pairs produce distinct signatures; each
has nine positions on its diagonal and two directions, that is,
\(18\) preimages. A different number of signatures or a class of size
different from \(18\) would refute the classification.
\end{proof}

\begin{lemma}[Exact classification of the multiplicative channel]
\label{pf84:E02-19}
\[
 |\operatorname{im}e|=26.
\]
There are \(24\) classes of size \(8\), one of size \(24\), and one of size
\(108\).
\end{lemma}
\begin{proof}
Enumerating the \(9^2\cdot4=324\) seeds using the explicit table of \(\ell_\times\) gives \(26\) sextets
\((E_1,\ldots,E_6)\) and the histogram
\[
 24\cdot8+24+108=324.
\]
An additional, absent, or differently sized class would refute the census.
\end{proof}

\begin{theorem}[Exact injective factorization of the catalog]
\label{pf84:E02-20}
\[
|\operatorname{im}s|=18,\qquad
|\operatorname{im}e|=26,\qquad
|\operatorname{im}G|=468=18\cdot26.
\]
Each additive signature has 18 preimages. The multiplicative signatures possess
the histogram
\[
24\text{ size classes }8,\qquad
1\text{ of size }24,\qquad
1\text{ of size }108.
\]
\end{theorem}
\begin{proof}
The two preceding lemmas give the cardinalities of the signatures. In each emitted
component, the hundreds digit is \(\bar S_r\), while the tens digit is
\(\bar S_r+\bar E_r\bmod10\). Therefore,
\[
 \bar S_r=\left\lfloor\frac{U_r}{100}\right\rfloor,\qquad
 \bar E_r=
 \left(\left\lfloor\frac{U_r}{10}\right\rfloor-\bar S_r\right)\bmod10.
\]
The units digit verifies \(3S_r+5E_r+1\bmod10\); it introduces no other
identification. Thus, \(G\) is injective, and its image has
\(18\cdot26=468\) elements. The \(24\) multiplicative classes of size
\(8\), combined with the \(18\) additive signatures, produce \(432\) sextets
of multiplicity \(144\); the other two classes produce \(18\) sextets of
multiplicity \(432\) and \(18\) of multiplicity \(1944\). In particular,
\[
 432\cdot144+18\cdot432+18\cdot1944=104\,976.
\]
Two pairs of signatures with the same \(U_6\), or a histogram that
failed to recover the complete domain, would refute the theorem.
\end{proof}

The complete composition factors, therefore, as
\[
\mathcal S^2\xrightarrow{s\times e}
\operatorname{im}s\times\operatorname{im}e
\xrightarrow{G}\mathcal U_6
\xrightarrow{\rho_3}\mathcal W_6\subset\mathbb F_3^6.
\]
The final projection has 243 values. The first descent preserves the
sum--product separation; the second preserves 468 decimal states; the
third removes the fiber coordinate and retains a ternary word.

\begin{corollary}[Census of the observable projection]
\label{cor:censo-proyeccion-observable}
\label{pf84:E02-21}
The complete enumeration determines the chain of cardinalities
\[
 104\,976\longrightarrow468\longrightarrow243\subset729.
\]
The \(468\) intermediate elements are sextuples of integers between
zero and \(999\); only their final image belongs to \(\mathbb F_3^6\). The
fiber multiplicities of the first map satisfy
\[
 432\cdot144+18\cdot432+18\cdot1944=104\,976.
\]
\end{corollary}

\begin{proof}
The factorization theorem gives \(18\cdot26=468\), and the exhaustive enumeration
of the \(104\,976\) seeds produces exactly the three
indicated multiplicities. Componentwise reduction traverses the
\(243\) words recorded in the ternary catalogue. The appended catalogues of
signatures and routes reproduce both enumerations.
\end{proof}

\begin{proposition}[Quintuple microfiber of the word \(010211\)]
\label{pf84:E02-22}
For the restricted reduction
\[
 \Psi:\mathcal U_6^{\rm cat}\longrightarrow\mathbb F_3^6,
 \qquad \Psi(U_6)=(-U_6)\bmod3,
\]
the fiber of \(010211\) consists of exactly the following five sextets:
\begingroup
\scriptsize
\[
\begin{array}{@{}cclc@{}}
\toprule
U_6&\sigma_+&\sigma_\times&\text{multiplicity}\\
\midrule
(501,614,498,169,272,272)&(5,6,4,1,2,2)&(5,5,5,5,5,5)&432\\
(810,923,870,169,272,272)&(8,9,8,1,2,2)&(3,3,9,5,5,5)&144\\
(810,983,810,169,272,272)&(8,9,8,1,2,2)&(3,9,3,5,5,5)&144\\
(870,923,810,169,272,272)&(8,9,8,1,2,2)&(9,3,3,5,5,5)&144\\
(870,923,810,418,674,521)&(8,9,8,4,6,5)&(9,3,3,7,1,7)&144\\
\bottomrule
\end{array}
\]
\endgroup
They are five preimages of the catalog \(U_6\), not five complete TPK
histories. They preserve distinct signatures and multiplicities although they share the same
visible word.
\end{proposition}

\begin{proof}
The exhaustive scan of the catalogue of \(468\) sextuples selects exactly
the five rows of the table. The signature-recovery function reproduces
the two sextuples of every row, and the sum of multiplicities is
\(432+4\cdot144=1008\). Enumeration additionally verifies that every cell is the
Cartesian product of its channel seeds. A cardinality other than five
or identification of two rows despite their enriched data would refute the
statement.
\end{proof}
